\documentclass[12pt,a4paper]{article}
\usepackage[margin=1in]{geometry}
\usepackage{graphicx}
\usepackage{listings}
\usepackage{xcolor}
\usepackage{amsmath}
\usepackage{amssymb}
\usepackage{hyperref}
\usepackage{booktabs}
\usepackage{float}

% Code styling
\lstset{
    language=C,
    basicstyle=\ttfamily\small,
    keywordstyle=\color{blue},
    commentstyle=\color{green!60!black},
    stringstyle=\color{red},
    numbers=left,
    numberstyle=\tiny,
    frame=single,
    breaklines=true,
    showstringspaces=false
}
% ---- TRIMMED VERSION ----

\begin{document}

\begin{titlepage}
    \centering
    \vspace*{2cm}

    {\LARGE \textbf{Lab 7: Acquire and Filter Analog Signals Using ADC and CMSIS DSP}}\par
    \vspace{1.5cm}

    {\large Name: Azyan Ahmed}\par
    \vspace{0.5cm}
    {\large ID: aa10287}\par
    \vspace{0.5cm}
    {\large \today}\par

    \vfill
\end{titlepage}

% ============================================================
\section{Objective}
% ============================================================
The objective of this lab is to understand and apply Analog-to-Digital Conversion (ADC) on the STM32F3 Discovery board and to implement digital signal filtering using the CMSIS DSP library. Specifically, we aim to:
\begin{itemize}
    \item Understand how ADC works and configure it on the STM32F303VCT6.
    \item Read and interpret analog values from a potentiometer sensor.
    \item Transmit sensor data over UART for real-time monitoring.
    \item Implement a 10-point Moving Average Low-Pass Filter using the CMSIS DSP library.
    \item Acquire a sine wave via ADC and visualize raw vs.\ filtered signals using a Python plotting script.
    \item Implement and compare polling-based and interrupt-based ADC acquisition.
\end{itemize}

% ============================================================
\section{Background: Analog-to-Digital Conversion}
% ============================================================

An Analog-to-Digital Converter (ADC) transforms a continuous analog input voltage into a discrete digital value that a microcontroller can process.

\subsection{Sampling and Quantization}
\textbf{Sampling} measures the analog signal amplitude at discrete time intervals. According to the \textbf{Nyquist Theorem}, the sampling rate must be at least twice the maximum signal frequency:
\begin{equation}
    f_s \geq 2 \cdot f_{max}
\end{equation}

\textbf{Quantization} maps each sample to the nearest of $2^n$ discrete levels for an $n$-bit ADC. The step size (LSB) and digital output are:
\begin{equation}
    \Delta V = \frac{V_{ref}}{2^n}
    \qquad
    D = \left\lfloor \frac{V_{in}}{V_{ref}} \cdot (2^n - 1) \right\rfloor
\end{equation}

\subsection{Key ADC Parameters}

\begin{table}[H]
\centering
\begin{tabular}{ll}
\toprule
\textbf{Parameter} & \textbf{Description} \\
\midrule
Resolution         & Number of bits; determines $2^n$ discrete levels \\
Reference Voltage ($V_{ref}$)  & Maximum voltage corresponding to full-scale output \\
Step Size ($\Delta V$) & Smallest detectable voltage change: $V_{ref}/2^n$ \\
Sampling Time       & Duration to sample before converting; longer = more stable \\
Data Alignment     & Right- or left-aligned result in the data register \\
Continuous Mode    & Auto-restarts conversion after each completed read \\
\bottomrule
\end{tabular}
\caption{Key ADC Configuration Parameters}
\end{table}

% ============================================================
\section{System Configuration}
% ============================================================

\subsection{ADC1 Configuration}

ADC1 is configured on Channel~1 (pin PA0, single-ended) with the following CubeMX settings:

\begin{table}[H]
\centering
\begin{tabular}{ll}
\toprule
\textbf{Parameter} & \textbf{Value} \\
\midrule
Resolution             & 12-bit (0--4095) \\
Data Alignment         & Right-aligned \\
Continuous Conv.\ Mode  & Enabled \\
EOC Selection          & End of single channel conversion \\
Sampling Time          & 601.5 cycles \\
External Trigger       & None (software start) \\
DMA                    & Disabled \\
NVIC Interrupt         & ADC1 and ADC2 global interrupt enabled \\
\bottomrule
\end{tabular}
\caption{ADC1 Configuration Parameters}
\end{table}

With 12-bit resolution and $V_{ref} = 3.3$\,V, the step size is $\Delta V = 3.3/4096 \approx 0.806$\,mV per LSB. The input voltage is recovered from raw reading $D$ by:
\begin{equation}
    V_{in} = D \times \frac{3.3}{4095}
\end{equation}

\subsection{USART2 Configuration}

\begin{table}[H]
\centering
\begin{tabular}{ll}
\toprule
\textbf{Parameter} & \textbf{Value} \\
\midrule
Mode            & Asynchronous \\
Baud Rate       & 115200 \\
Word Length      & 8 bits \\
Stop Bits       & 1 \\
Parity          & None \\
\bottomrule
\end{tabular}
\caption{USART2 Configuration Parameters}
\end{table}


% ============================================================
\section{Task 1: Reading an Analog Sensor and Transmitting via UART}
% ============================================================

\subsection{Objective}
Read the analog voltage from a potentiometer connected to ADC1 Channel~1 (PA0) and transmit the raw digital value over UART to a serial terminal in real time.

\subsection{Hardware Setup}
A potentiometer is wired as a voltage divider:
\begin{itemize}
    \item One side pin $\rightarrow$ 3.3\,V
    \item Other side pin $\rightarrow$ GND
    \item Middle pin (wiper) $\rightarrow$ PA0 (ADC1\_IN1)
\end{itemize}

As the knob is rotated, the wiper voltage varies continuously between 0\,V and 3.3\,V, which the ADC converts to a value between 0 and 4095.

\subsection{Approach}
\begin{enumerate}
    \item Calibrate ADC1 in single-ended mode using \texttt{HAL\_ADCEx\_Calibration\_Start()}.
    \item Start the ADC in polling mode with \texttt{HAL\_ADC\_Start()}.
    \item In the main loop, poll for conversion completion, read the value, format it as a string, and transmit over UART.
    \item Add a 10\,ms delay to limit the output to approximately 100 readings per second.
\end{enumerate}

\subsection{Code Implementation}

\begin{lstlisting}[caption=Task 1: Includes and Initialisation]
/* USER CODE BEGIN Includes */
#include <stdint.h>
#include <stdio.h>
#include <string.h>
/* USER CODE END Includes */

/* USER CODE BEGIN 2 */
// Calibrate ADC1 in single-ended mode
HAL_ADCEx_Calibration_Start(&hadc1, ADC_SINGLE_ENDED);
// Start ADC1 in polling mode
HAL_ADC_Start(&hadc1);
/* USER CODE END 2 */
\end{lstlisting}

\begin{lstlisting}[caption=Task 1: Main Loop -- Polling ADC and UART Transmission]
while (1)
{
    // Wait for ADC1 to finish conversion
    HAL_ADC_PollForConversion(&hadc1, HAL_MAX_DELAY);

    // Read ADC1 converted value (0-4095)
    uint32_t adc_ki_value = HAL_ADC_GetValue(&hadc1);

    char adc_str[50];

    // Format the value into a readable string
    int len = snprintf(adc_str, sizeof(adc_str),
                       "ADC Value: %lu\r\n",
                       (unsigned long)adc_ki_value);

    // Transmit over UART2
    HAL_UART_Transmit(&huart2, (uint8_t*)adc_str,
                       len, HAL_MAX_DELAY);

    // Delay 10 ms (~100 samples/second)
    HAL_Delay(10);
}
\end{lstlisting}

\textbf{Key functions:} \texttt{HAL\_ADCEx\_Calibration\_Start()} corrects internal ADC offset errors before use; \texttt{HAL\_ADC\_PollForConversion()} blocks until the EOC flag is set; \texttt{HAL\_ADC\_GetValue()} reads the 12-bit result; \texttt{HAL\_UART\_Transmit()} sends the formatted string over USART2.

\subsection{Results}
The ADC successfully read values from the potentiometer across the full 12-bit range. At one extreme of the potentiometer rotation, the ADC output was approximately 0 (corresponding to 0\,V). At the other extreme, the output reached approximately 4095 (corresponding to 3.3\,V). Intermediate positions produced proportional values, confirming correct linear operation of the voltage divider and ADC.

\begin{figure}[H]
    \centering
    \includegraphics[width=0.6\linewidth]{task1_output.png}
    \caption{Task 1: UART terminal output showing ADC values changing as the potentiometer is rotated}
\end{figure}

% ============================================================
\section{Task 2: Signal Filtering Using CMSIS DSP Library}
% ============================================================

\subsection{Objective}
Apply a 10-point Moving Average Low-Pass Filter to a sine wave input acquired via ADC, and visualize both the raw and filtered signals in real time using a Python plotting script. The task is completed in two parts: polling mode and interrupt mode.

\subsection{Introduction to Filtering}
Filtering in signal processing modifies or removes unwanted components from a signal. Common filter types include:
\begin{itemize}
    \item \textbf{Low-Pass Filter (LPF):} Passes low-frequency components, attenuates high-frequency components. Used to smooth noisy signals.
    \item \textbf{High-Pass Filter (HPF):} Passes high-frequency components, blocks low frequencies.
    \item \textbf{Band-Pass Filter:} Passes a specific frequency range.
    \item \textbf{Moving Average Filter:} A simple non-recursive LPF that computes the average of the most recent $N$ samples.
\end{itemize}

\subsection{Moving Average Filter Theory}

A moving average filter of length $N$ computes the output $y[n]$ as:
\begin{equation}
    y[n] = \frac{1}{N} \sum_{k=0}^{N-1} x[n-k]
\end{equation}

where $x[n-k]$ represents the most recent $N$ input samples. In this lab, $N = 10$.

\subsection{CMSIS DSP Library}
The Cortex Microcontroller Software Interface Standard (CMSIS) provides optimized DSP functions for ARM Cortex-M processors. We use \texttt{arm\_mean\_f32()} from the statistics module, which computes the arithmetic mean of an array of \texttt{float32\_t} values using hardware FPU acceleration on the Cortex-M4F.

\subsection{Hardware Setup}
\begin{itemize}
    \item A function generator is connected to PA0 (ADC1 Channel~1).
    \item The function generator outputs a sine wave with frequency varied between 1\,Hz and 100\,Hz.
    \item The sine wave amplitude is kept within the ADC input range (0--3.3\,V), typically with a 1.65\,V DC offset and 1.5\,V peak-to-peak amplitude.
    \item The STM32 board is connected to the PC via USB for both flashing and UART data transmission.
\end{itemize}

\subsection{Variable Declarations}

\begin{lstlisting}[caption=Task 2: Global Variables and Filter Definition]
/* USER CODE BEGIN PD */
#define FILTER_LEN 10
/* USER CODE END PD */

/* USER CODE BEGIN PV */
// Circular buffer for the last 10 ADC samples
float32_t inputBuffer[FILTER_LEN];
float32_t filtered_output;
uint8_t bufferIndex = 0;
/* USER CODE END PV */
\end{lstlisting}

\subsection{Moving Average Filter Function}

\begin{lstlisting}[caption=Task 2: Moving Average Filter Implementation]
/* USER CODE BEGIN 0 */
void apply_moving_average(float32_t new_sample) {
    // Store the new sample at the current index
    inputBuffer[bufferIndex] = new_sample;

    // Advance index circularly
    bufferIndex = (bufferIndex + 1) % FILTER_LEN;

    // Compute the mean of the buffer using CMSIS DSP
    arm_mean_f32(inputBuffer, FILTER_LEN, &filtered_output);
}
/* USER CODE END 0 */
\end{lstlisting}

\textbf{Explanation:}
\begin{itemize}
    \item \texttt{inputBuffer[]} acts as a \textbf{circular buffer}: the oldest sample is overwritten each time a new sample arrives, and \texttt{bufferIndex} wraps around using the modulo operator.
    \item \texttt{arm\_mean\_f32(inputBuffer, FILTER\_LEN, \&filtered\_output)} computes the arithmetic mean of all 10 values in the buffer and stores the result in \texttt{filtered\_output}. This is mathematically equivalent to the moving average formula.
    \item The circular buffer ensures that only the most recent $N = 10$ samples contribute to the average at any given time.
\end{itemize}

% ────────────────────────────────────────────────────────────
\subsection{Part (i): Polling Mode Implementation}
% ────────────────────────────────────────────────────────────

\begin{lstlisting}[caption=Task 2 -- Polling Mode: Initialisation]
/* USER CODE BEGIN 2 */
HAL_ADCEx_Calibration_Start(&hadc1, ADC_SINGLE_ENDED);
HAL_ADC_Start(&hadc1);
/* USER CODE END 2 */
\end{lstlisting}

\begin{lstlisting}[caption=Task 2 -- Polling Mode: Main Loop]
while (1)
{
    // Wait for conversion to complete
    HAL_ADC_PollForConversion(&hadc1, HAL_MAX_DELAY);
    uint32_t adc_raw = HAL_ADC_GetValue(&hadc1);

    // Convert raw ADC reading to voltage (0-3.3V)
    float32_t voltage = (float32_t)adc_raw * (3.3f / 4095.0f);

    // Apply 10-point moving average filter
    apply_moving_average(voltage);

    // Convert filtered voltage back to ADC scale for plotting
    uint32_t filtered_int =
        (uint32_t)(filtered_output * (4095.0f / 3.3f));

    // Transmit CSV format: rawValue,filteredValue
    char msg[64];
    int len = snprintf(msg, sizeof(msg), "%lu,%lu\r\n",
                       (unsigned long)adc_raw,
                       (unsigned long)filtered_int);
    HAL_UART_Transmit(&huart2, (uint8_t*)msg, len, HAL_MAX_DELAY);

    HAL_Delay(10);
}
\end{lstlisting}

\textbf{Data flow:}
\begin{enumerate}
    \item The raw ADC value (0--4095) is read via polling.
    \item It is converted to a floating-point voltage: $V = D \times (3.3 / 4095)$.
    \item The voltage is passed through the 10-point moving average filter.
    \item The filtered voltage is converted back to the 0--4095 ADC scale for consistent plotting.
    \item Both raw and filtered values are transmitted as comma-separated values (CSV) over UART.
\end{enumerate}

% ────────────────────────────────────────────────────────────
\subsection{Part (ii): Interrupt Mode Implementation}
% ────────────────────────────────────────────────────────────

In interrupt mode, the ADC is started with \texttt{HAL\_ADC\_Start\_IT()} and the conversion-complete callback handles all processing asynchronously.

\begin{lstlisting}[caption=Task 2 -- Interrupt Mode: Initialisation]
/* USER CODE BEGIN 2 */
HAL_ADCEx_Calibration_Start(&hadc1, ADC_SINGLE_ENDED);
HAL_ADC_Start_IT(&hadc1);
/* USER CODE END 2 */
\end{lstlisting}

\begin{lstlisting}[caption=Task 2 -- Interrupt Mode: Main Loop]
while (1)
{
    // CPU is free; all ADC work is done in the callback
    HAL_Delay(10);
}
\end{lstlisting}

\begin{lstlisting}[caption=Task 2 -- Interrupt Mode: ADC Callback]
/* USER CODE BEGIN 4 */
void HAL_ADC_ConvCpltCallback(ADC_HandleTypeDef *hadc) {
    if (hadc->Instance == ADC1) {
        // Get the raw ADC value
        uint32_t adc_raw = HAL_ADC_GetValue(hadc);

        // Convert to voltage (0-3.3V scale)
        float32_t voltage =
            (float32_t)adc_raw * (3.3f / 4095.0f);

        // Apply 10-point moving average filter
        apply_moving_average(voltage);

        // Convert filtered voltage back to ADC scale
        uint32_t filtered_int =
            (uint32_t)(filtered_output * (4095.0f / 3.3f));

        // Transmit: rawValue,filteredValue
        char msg[64];
        int len = snprintf(msg, sizeof(msg), "%lu,%lu\r\n",
                           (unsigned long)adc_raw,
                           (unsigned long)filtered_int);
        HAL_UART_Transmit(&huart2, (uint8_t*)msg,
                           len, HAL_MAX_DELAY);
    }
}
/* USER CODE END 4 */
\end{lstlisting}

\textbf{Key differences from polling mode:}
\begin{itemize}
    \item \texttt{HAL\_ADC\_Start\_IT()} is used instead of \texttt{HAL\_ADC\_Start()}, which enables the ADC end-of-conversion interrupt.
    \item The interrupt vector \texttt{ADC1\_2\_IRQHandler()} in \texttt{stm32f3xx\_it.c} calls \texttt{HAL\_ADC\_IRQHandler(\&hadc1)}, which internally invokes our \texttt{HAL\_ADC\_ConvCpltCallback()} when the EOC flag is set.
    \item The main loop is nearly empty---all ADC reading, filtering, and UART transmission occur inside the ISR callback.
    \item The CPU is free to perform other tasks between conversions, making this approach more efficient for real-time systems.
\end{itemize}

% ────────────────────────────────────────────────────────────
\subsection{Python Visualization Script}
% ────────────────────────────────────────────────────────────

A Python script (\texttt{plot\_adc.py}) reads the CSV-formatted UART data in real time and plots both signals using \texttt{matplotlib}:

\begin{lstlisting}[language=Python,caption=Python Plotting Script (plot\_adc.py)]
import serial
import matplotlib.pyplot as plt
from drawnow import *

# Setup Serial
sinWaveData = serial.Serial('/dev/ttyUSB0', 115200)
plt.ion()  # Enable interactive mode

adcValues = []
filteredValues = []
time_ms = []
cnt = 0
SAMPLE_PERIOD_MS = 10  # matches HAL_Delay(10)

def makeFig():
    plt.clf()
    plt.title('Live ADC & Filtered Data')
    plt.grid(True)
    plt.xlabel('Time (ms)')
    plt.ylabel('ADC Value')
    plt.ylim(0, 5000)
    plt.plot(time_ms, adcValues, 'r.-', label='Raw ADC')
    plt.plot(time_ms, filteredValues, 'b.-',
             label='Filtered Mean')
    plt.legend(loc='upper left')

while True:
    while sinWaveData.inWaiting() == 0:
        pass  # Wait for data

    try:
        line = sinWaveData.readline().decode().strip()
        values = line.split(',')

        if len(values) == 2:
            raw = int(values[0])
            filtered = int(values[1])

            adcValues.append(raw)
            filteredValues.append(filtered)
            time_ms.append(cnt * SAMPLE_PERIOD_MS)
            cnt += 1

            drawnow(makeFig)

            # Keep only last 500 samples
            if len(adcValues) > 500:
                adcValues.pop(0)
                filteredValues.pop(0)
                time_ms.pop(0)
    except Exception as e:
        print("Error:", e)
\end{lstlisting}

\subsection{Results}

\subsubsection{Polling Mode Output}
In polling mode, the raw and filtered ADC values were transmitted over UART and plotted in real time. The raw signal showed the expected sine waveform with minor noise, while the filtered signal appeared smoother with reduced high-frequency fluctuations.

\begin{figure}[H]
    \centering
    \includegraphics[width=0.7\linewidth]{task2_polling.png}
    \caption{Task 2 -- Polling Mode: Raw (red) and filtered (blue) ADC signals}
\end{figure}

\subsubsection{Interrupt Mode Output}
In interrupt mode, identical results were obtained with the same quality of filtering. The key difference was in CPU utilisation: the main loop was effectively idle, with all processing handled in the ISR callback.

\begin{figure}[H]
    \centering
    \includegraphics[width=0.7\linewidth]{task2_interrupt.png}
    \caption{Task 2 -- Interrupt Mode: Raw (red) and filtered (blue) ADC signals}
\end{figure}

\subsubsection{Effect of Increasing Frequency}
When the input sine wave frequency was progressively increased on the function generator, the amplitude of the filtered output signal decreased. This is the expected behaviour of a low-pass filter: higher-frequency components are attenuated more strongly.

\begin{figure}[H]
    \centering
    \includegraphics[width=0.7\linewidth]{task2_freq_sweep.png}
    \caption{Filtered output amplitude decreasing as input frequency increases, demonstrating low-pass filter behaviour}
\end{figure}

\subsection{Polling vs.\ Interrupt Mode Comparison}

\begin{table}[H]
\centering
\begin{tabular}{lcc}
\toprule
\textbf{Aspect} & \textbf{Polling} & \textbf{Interrupt} \\
\midrule
CPU utilisation during conversion  & Blocked (100\%) & Free ($\sim$0\%) \\
Code complexity                     & Simpler & Slightly more complex \\
Suitable for multi-tasking          & No  & Yes \\
Deterministic timing                & Yes (blocking) & Depends on ISR priority \\
Implementation location             & Main loop & Callback function \\
Function used                       & \texttt{HAL\_ADC\_Start()} & \texttt{HAL\_ADC\_Start\_IT()} \\
\bottomrule
\end{tabular}
\caption{Comparison of Polling and Interrupt ADC Modes}
\end{table}

% ============================================================
\section{Evaluation Questions}
% ============================================================

\subsection*{(a) Explain the relationship between ADC resolution and quantization error. Why does increasing resolution reduce error?}

\textbf{Relationship:}

Quantization error arises because the ADC maps a continuous range of analog voltages to a finite set of discrete digital levels. Each analog voltage is rounded to the nearest quantization level, and the maximum rounding error is $\pm \frac{1}{2}$ LSB.

For an $n$-bit ADC with reference voltage $V_{ref}$, the step size (LSB) is:
\begin{equation}
    \Delta V = \frac{V_{ref}}{2^n}
\end{equation}

The maximum quantization error is therefore:
\begin{equation}
    \epsilon_{max} = \pm \frac{\Delta V}{2} = \pm \frac{V_{ref}}{2^{n+1}}
\end{equation}

\textbf{Numerical Example:}

For $V_{ref} = 3.3$\,V:
\begin{itemize}
    \item \textbf{8-bit ADC:} $\Delta V = 3.3 / 256 = 12.89$\,mV, $\epsilon_{max} = \pm 6.45$\,mV
    \item \textbf{10-bit ADC:} $\Delta V = 3.3 / 1024 = 3.22$\,mV, $\epsilon_{max} = \pm 1.61$\,mV
    \item \textbf{12-bit ADC:} $\Delta V = 3.3 / 4096 = 0.806$\,mV, $\epsilon_{max} = \pm 0.403$\,mV
\end{itemize}

\textbf{Why increasing resolution reduces error:}

Each additional bit of resolution \emph{doubles} the number of quantization levels ($2^n$), which \emph{halves} the step size and therefore halves the maximum quantization error. In other words, the step size decreases \emph{exponentially} with the number of bits:
\begin{equation}
    \Delta V \propto 2^{-n}
\end{equation}

Going from 8~bits to 12~bits increases the number of levels $16\times$ (from 256 to 4096), reducing the quantization error by the same factor. The signal can be represented with much finer granularity, meaning the digital approximation is much closer to the true analog value.

In the context of this lab, using 12-bit resolution gives us a step size of approximately 0.8\,mV, which is more than sufficient to accurately represent the potentiometer voltage and capture the subtle variations in a filtered sine wave signal.

% ============================================================
\section{Conclusion}
% ============================================================

This lab demonstrated two core embedded systems skills: analog-to-digital conversion and real-time digital signal filtering on the STM32F3 Discovery board.

\textbf{Task 1} established the foundation by reading a potentiometer via ADC1 in polling mode and transmitting the raw 12-bit values over UART. The full range (0--4095) was confirmed across the potentiometer travel, with a quantization step of $\approx 0.8$\,mV at $V_{ref} = 3.3$\,V.

\textbf{Task 2} applied a 10-point moving average low-pass filter using \texttt{arm\_mean\_f32()} from the CMSIS DSP library. The filter was implemented in both polling and interrupt modes. Interrupt mode freed the CPU during conversions, with all ADC reading, filtering, and UART transmission handled entirely inside \texttt{HAL\_ADC\_ConvCpltCallback()}. When the input sine wave frequency was increased, the filtered output amplitude progressively decreased, confirming the low-pass behaviour of the moving average filter.

Real-time visualisation was achieved using a Python script (\texttt{plot\_adc.py}) with \texttt{pyserial} and \texttt{matplotlib}, plotting both raw and filtered signals simultaneously.

\end{document}
